% Figure for Calculus §57.
\begin{tikzpicture}
  \begin{axis}[nfig, width=8.6cm, height=7cm, xmin=-4, xmax=4, ymin=-4, ymax=4,
      axis lines=middle, xtick={-4,-2,2,4}, ytick={-4,-2,2,4},
      xlabel={$t$}, ylabel={$y$}, restrict y to domain=-6:6, unbounded coords=jump]
    \addplot[black!35, thin, dashed, domain=-4:4] {1};
    \addplot[black!35, thin, dashed, domain=-4:4] {-1};
    % members with c < 0: between -1 and 1
    \foreach \c in {-0.05,-0.4,-3} {
      \addplot[figB!70, domain=-4:4, samples=120] {(1+\c*exp(x))/(1-\c*exp(x))};
    }
    \addplot[figB!70, domain=-4:4] {1};
    % members with c > 0: blow up at t = -ln c
    \foreach \c in {0.04, 3} {
      \pgfmathsetmacro{\s}{-ln(\c)}
      \addplot[figB!70, domain=-4:\s-0.02, samples=150] {(1+\c*exp(x))/(1-\c*exp(x))};
      \addplot[figB!70, domain=\s+0.02:4, samples=150] {(1+\c*exp(x))/(1-\c*exp(x))};
    }
    % the IVP solution, c = 1/3
    \addplot[figR, thick, domain=-4:1.06, samples=200] {(3+exp(x))/(3-exp(x))};
    \addplot[figR, thick, domain=1.14:4, samples=200] {(3+exp(x))/(3-exp(x))};
    \draw[figR, thin, dotted] (axis cs:1.0986,-4) -- (axis cs:1.0986,4);
    \node[font=\scriptsize, anchor=south west] at (axis cs:2.6,1.0) {$y=1$};
    \node[black!60, font=\scriptsize, anchor=north west] at (axis cs:-3.95,-1.0) {$y=-1$};
    \fill[figR] (axis cs:0,2) circle (2pt);
    \node[figR, font=\scriptsize, anchor=east] at (axis cs:-0.1,2.2) {$(0,2)$};
    \node[figR, font=\scriptsize, anchor=west] at (axis cs:1.15,3.4) {$t = \ln 3$};
  \end{axis}
\end{tikzpicture}
